\documentclass[11pt]{article}
\usepackage{amsmath,amssymb,amsthm,hyperref}
\newtheorem{theorem}{Theorem}
\newtheorem{lemma}[theorem]{Lemma}
\title{An Arbitrarily Compact Universal Copula Module}
\author{}
\date{}
\begin{document}
\maketitle

\paragraph{Status.}
This is a finite theorem with a complete written proof and exact symbolic
verification in \texttt{audit\_copula\_compact\_limit.py}. It is not
Lean-checked and does not prove an asymptotic saving for unrestricted
head complexity. It strengthens the finite module in
\texttt{continuation-copula-trades.tex}: all memory slopes can now tend
to zero while a fixed family of three denominator templates still
programs all eight mixed-fiber signs. The attention model and its
admissibility conditions are unchanged.

\begin{theorem}\label{thm:compact}
There is $\varepsilon_0>0$ such that, for every
$0<\varepsilon<\varepsilon_0$ and every prescribed sign at each of
the eight vertices
\[
 (u,v_1,v_2,v_3)\in\{0,1\}^4,
 \qquad v_1+v_2+v_3\ge2,
\]
four original admissible heads realize those signs strictly and are
exactly zero at all eight remaining vertices. Their denominators are
\[
 1-\frac{u}{16}-\varepsilon
       \sum_{j=1}^3(\pi_j(h)+1)v_j,\qquad h=0,1,2,3,
 \tag{1}
\]
where $\pi_1=\mathrm{id}$ and one of these three pairs is chosen:
\[
\begin{array}{c|c|c}
 t&\pi_2&\pi_3\\\hline
 0&(0,1,2,3)&(0,1,2,3)\\
 1&(0,2,1,3)&(1,0,3,2)\\
 2&(0,1,3,2)&(1,2,3,0).
\end{array}
\tag{2}
\]
The numerators have only $v_j$ terms. For rational $\varepsilon$ all
parameters can be chosen rational. We may take $\varepsilon_0\le1/32$,
so every denominator is at least $9/16$ on the whole Boolean cube.
\end{theorem}

\paragraph{Residue space and exact zero fibers.}
Put $d_0=1$, $d_1=15/16$, and, for $r=1,2,3,4$ and $k=0,1$, set
\[
 q_k(r;\varepsilon)=
 \frac{(\varepsilon r-d_0)(\varepsilon r-d_1)r^k}
      {\prod_{\ell\in\{1,2,3,4\}\setminus\{r\}}(r-\ell)}.
 \tag{3}
\]
For either $d=d_0,d_1$,
\[
 \sum_{r=1}^4\frac{q_k(r;\varepsilon)}{d-\varepsilon r}=0.
 \tag{4}
\]
To verify (4), cancel the factor $\varepsilon r-d$ in (3).
The remaining numerator is a polynomial in $r$ of degree at most two.
For every such polynomial $p$,
$\sum_r p(r)/\prod_{\ell\ne r}(r-\ell)=0$: it is the coefficient
of $X^3$ in the Lagrange interpolation formula for $p(X)$.
The vectors $q_0,q_1$ are independent for sufficiently small
$\varepsilon$, since their entrywise ratio is $r$ and all entries of
$q_0$ are nonzero.

Choose six free real numbers $\theta_{jk}$ and put
\[
 c_j(r)=\sum_{k=0}^1\theta_{jk}q_k(r;\varepsilon),\qquad
 F(u,v)=\sum_{h=0}^3
 \frac{\sum_{j=1}^3c_j(\pi_j(h)+1)v_j}
 {1-u/16-\varepsilon\sum_{j=1}^3(\pi_j(h)+1)v_j}.
 \tag{5}
\]
Equation (4) proves exact vanishing when $v$ has one nonzero coordinate;
the numerator proves vanishing when $v=0$.

\paragraph{Polynomial matrices.}
Order the eight mixed outputs by masks $S=3,5,6,7$, with binary
weights $1,2,4$ for the memory coordinates, and within each mask by
$d=d_0,d_1$. For a fixed template, write
$T_h(S)=\sum_{j\in S}(\pi_j(h)+1)$. The output matrix is
\[
 A_{(S,d),(j,k)}(\varepsilon)=
 {\bf1}_{j\in S}\sum_{h=0}^3
 \frac{q_k(\pi_j(h)+1;\varepsilon)}{d-\varepsilon T_h(S)}.
 \tag{6}
\]
Multiplying row $(S,d)$ by the positive quantity
$Q_{S,d}=\prod_h(d-\varepsilon T_h(S))$ gives the polynomial matrix
\[
 \widetilde B_{(S,d),(j,k)}=
 {\bf1}_{j\in S}\sum_h q_k(\pi_j(h)+1;\varepsilon)
                 \prod_{g\ne h}(d-\varepsilon T_g(S)).
 \tag{7}
\]
Divide each column by the polynomial greatest common divisor of its
eight entries to obtain $B$. Every divisor is a nonzero polynomial,
so for all sufficiently small positive $\varepsilon$ the images of
$B$ and $A$ meet exactly the same strict orthants.

\paragraph{Exact limiting kernel computation.}
For each of the three templates, $B$ has rank six over
$\mathbb Q(\varepsilon)$. A polynomial basis $K(\varepsilon)$ for
its two-dimensional left kernel can be normalized so that $K(0)$
has rank two. The following matrices have precisely the row spaces
of these three $K(0)$:
\[
 L_0=\begin{pmatrix}
 43065&0&43065&0&43065&0&-124740&103936\\
 0&45936&0&45936&0&45936&-126225&102544
 \end{pmatrix},
 \tag{8}
\]
\[
 L_1=\begin{pmatrix}
 200160&0&15076035&-16604160&-14106240&18333696&-5866560&5161984\\
 0&17792&1285875&-1415808&-1231200&1607040&-529200&466560
 \end{pmatrix},
 \tag{9}
\]
\[
 L_2=\begin{pmatrix}
 3345&0&-57855&69632&31245&-31744&-13920&12032\\
 0&3568&-61200&73200&27900&-28176&-10575&8464
 \end{pmatrix}.
 \tag{10}
\]
Here is an explicit reproducible computation of the limits, including
how cancellations at $\varepsilon=0$ are handled. Choose the
lexicographically first six-row set $R$ with
$\Delta=\det B_R\ne0$. For each row $s\notin R$, make the vector
$z^{(s)}$ whose $s$ entry is $\Delta$, whose other entry outside
$R$ is zero, and whose entry at the $i$th row of $R$ is
\[
 -\det B_{R[i\leftarrow s]}.
 \tag{11}
\]
The determinant in (11) retains the displayed row order. Cramer's
rule gives $B^\top z^{(s)}=0$. Divide each vector by its polynomial
gcd, obtaining $a,b$. Choose the first index $p$ with $a_p(0)\ne0$;
replace $b$ by $a_pb-b_pa$ and again divide by its polynomial gcd.
The resulting two polynomial vectors are $K$'s rows. Evaluating at
zero, taking reduced row echelon form, clearing rational denominators
and dividing each row by its integer gcd gives (8)--(10).
The companion checker performs exactly these operations over
$\mathbb Q[\varepsilon]$, verifies $KB=0$ by expansion, verifies
the nonzero minor, and checks the stated limiting rows.

\paragraph{Strict coverage by the limiting planes.}
The three six-dimensional spaces $\ker L_t$ jointly meet all 256
strict orthants of $\mathbb R^8$. The following exact procedure
verifies this assertion without optimization. For each sign word
$s\in\{-1,1\}^8$, and each $L=L_t$, set
\[
 M=L\operatorname{diag}(s),\qquad b=-M\mathbf1.
 \tag{12}
\]
Examine every pair of columns $M_i,M_j$ with nonzero determinant.
If
\[
 (\alpha,\beta)^\top=[M_i\ M_j]^{-1}b
 \quad\hbox{has }\alpha,\beta\ge0,
 \tag{13}
\]
then $z=\mathbf1+\alpha e_i+\beta e_j$ is strictly positive and
$L\operatorname{diag}(s)z=0$. Checking (13) in template order
$0,1,2$ accepts respectively $224,26,6$ new words, hence all 256.
All arithmetic in this check is rational, all output margins are at
least one, and the script can export every witness. The positive
vectors themselves suffice; no completeness claim about the pair
search is needed for this verification.

\paragraph{Passage to small positive memory slopes.}
Fix one sign word and one limiting witness $x\in\ker K(0)$ with
$s_ix_i\ge1$. Since $K(0)$ has rank two, the matrix
$K(\varepsilon)K(\varepsilon)^\top$ is invertible for sufficiently
small $\varepsilon$. Define the continuous projection
\[
 P(\varepsilon)=I-K(\varepsilon)^\top
 [K(\varepsilon)K(\varepsilon)^\top]^{-1}K(\varepsilon).
 \tag{14}
\]
Then $P(\varepsilon)x\to x$ and consequently
$s_i(P(\varepsilon)x)_i>0$ for every $i$ when $\varepsilon$ is
sufficiently small. Also
$P(\varepsilon)x\in\ker K(\varepsilon)=\operatorname{im}B(\varepsilon)$,
because the nonzero polynomial minor proves rank six away from
finitely many roots. Positive row rescaling in (7) now gives the
desired strict signs in (6). There are only 256 words, so the
minimum of their finitely many positive neighborhoods supplies one
$\varepsilon_0>0$ valid for all words. Taking it smaller than $1/32$
proves all the asserted denominator bounds. For rational
$\varepsilon$, (14) and a six-row linear solve are rational, so
all residues can be rational as well.

\paragraph{Original attention admissibility.}
For each summand of (5), write its numerator as $\sum_i t_ix_i$
and denominator as $1-\sum_i a_ix_i$, where
$a_u=1/16$ and $a_{v_j}=\varepsilon(\pi_j(h)+1)>0$.
Since $\varepsilon<1/32$, $\sum_i a_i<7/16$. The original atom
parameters
\[
 \alpha=\tfrac12,\quad\rho_i=2a_i,\quad
 \gamma=1-2\sum_i a_i,\quad\delta=0,\quad
 m_i=-t_i/a_i,\quad\eta=-\sum_i\rho_i m_i
 \tag{15}
\]
satisfy $\alpha>0$, $\rho_i>0$, and $\gamma>1/8$.
Since $\rho_i\alpha^{x_i}=2a_i-a_ix_i$ on Boolean inputs,
substitution into the original numerator and denominator gives
the summand of (5) exactly. Thus no zero weight, orientation
change, approximate admissibility, or additional head is used.

\paragraph{Remaining asymptotic issue.}
This theorem reuses four heads across four mixed memory fibers,
with two row evaluations each. Repetition of this fixed module
only changes a constant. An asymptotic improvement requires a
family whose ratio of independently programmed strict signs to
the available affine numerator parameters tends to infinity,
together with a whole-cube compiler at the corresponding head
cost. This theorem supplies neither assertion.
\end{document}
